%%% FIELD proof-boundary-frontier-obstruction
Put \(n=2m\).  By the balanced-dimension, even-benchmark, balanced-L-mask, and balanced-one-chart assumptions, the unique chart has
\[
C:=C_{b_0}=P:=P_{b_0}=\{1,\ldots,m\},\qquad
I:=I_{b_0}=Q:=Q_{b_0}=\{m+1,\ldots,2m\},
\]
where the repeated labels distinguish row from column sets, and
\[
W_n=I\times Q,qquad
\Omega_n=(C\times P)\mathbin{\sqcup}(I\times P)
\mathbin{\sqcup}(C\times Q),qquad |\Omega_n|=3m^2.       \tag{1}
\]
In particular, every cell used below belongs to the observed wedge whenever it lies in one of the three rectangles in (1).

We first identify the entire \(\mu=1\) boundary class.  By the model-class definition together with the known-rank, rank-one-benchmark, and strong-factor assumptions, every \(\vartheta_n=(L_n,D_n)\in\Theta_n\) admits a singular-value decomposition
\[
L_n=s_nu_nv_n^{\top},\qquad
\lVert u_n\rVert_2=\lVert v_n\rVert_2=1,qquad
s_-n\le s_n\le s_+n.                                    \tag{2}
\]
At \(\mu=1\), the incoherence assumption becomes
\[
|u_{n,i}|^2\le n^{-1}\quad(1\le i\le n),
\qquad |v_{n,t}|^2\le n^{-1}\quad(1\le t\le n).        \tag{3}
\]
The \(n\) nonnegative quantities on each side of (3) sum to one, so every inequality in (3) is an equality.  Consequently there are signs \(p_i,q_t\in\{-1,1\}\) and a level \(\ell_n=s_n/n\in[s_-,s_+]\) such that
\[
u_{n,i}=\frac{p_i}{\sqrt n},\qquad
v_{n,t}=\frac{q_t}{\sqrt n},\qquad
L_{n,it}=\ell_np_iq_t.                                  \tag{4}
\]
The two normalized overlap Grams in the unique chart are well defined because \(|C|=|P|=m>0\).  Equation (4) gives
\[
\frac{n}{|C|}u_{n,C}^{\top}u_{n,C}
=\frac{n}{m}\sum_{i\in C}\frac1n=1,
\qquad
\frac{n}{|P|}v_{n,P}^{\top}v_{n,P}
=\frac{n}{m}\sum_{t\in P}\frac1n=1.                  \tag{5}
\]
Thus every boundary matrix satisfies the overlap-Gram assumption for every \(0<\kappa\le1\), and the value of \(\kappa\) imposes no further restriction on this boundary class.

We next establish the claimed geometric lower bound.  The class is nonempty: choose \(p_i=q_t=1\), \(\ell_n=s_-\), and \(D_n=0\).  Equations (2)--(5), the treatment-support assumption, and \(0<\kappa\le1\) verify every member property in the model-class definition.  For this matrix, let \(q=(q_1,\ldots,q_n)^{\top}\), and take a vector \(h\in\mathbb R^n\) supported on \(I\).  Since \(v_n=q/\sqrt n\),
\[
\Delta=hq^{\top}=(\sqrt n\,h)v_n^{\top}\in\mathcal T_n(L_n)       \tag{6}
\]
by the tangent-space definition.  By (1) and the support of \(h\),
\[
\lVert P_{\Omega_n}\Delta\rVert_F
=\lVert h_Iq_P^{\top}\rVert_F
=\sqrt m\,\lVert h\rVert_2,
\qquad
\lVert P_{W_n}\Delta\rVert_*
=\lVert h_Iq_Q^{\top}\rVert_*
=\sqrt m\,\lVert h\rVert_2.                            \tag{7}
\]
Here a rank-one matrix \(ab^{\top}\) has its sole nonzero singular value equal to \(\lVert a\rVert_2\lVert b\rVert_2\), proving the nuclear-norm equality.  Choose \(h\) so that \(\sqrt m\lVert h\rVert_2=\sigma\sqrt{2n}\).  Since \(r=1\) and \(N_n=T_n=n\), this \(\Delta\) is feasible in the wedge-modulus definition.  The geometric-radius definition therefore yields
\[
\mathfrak R_n\ge \omega_n(L_n,\Omega_n)
\ge\sigma\sqrt{2n}.                                    \tag{8}
\]

For completeness, \(\mathfrak R_n\) is finite on this class.  For any boundary \(L_n\) and any \(\Delta\in\mathcal T_n(L_n)\), multiplication on the left and right by the diagonal sign matrices generated by \(p\) and \(q\) preserves Frobenius and nuclear norms and turns \(\Delta_{it}\) into an additive array \(A_{it}=a_i+d_t\).  Fix \(i_0\in C\) and \(t_0\in P\).  Then
\[
A_{it}=A_{i t_0}+A_{i_0t}-A_{i_0t_0}qquad(i\in I, t\in Q).       \tag{9}
\]
Using \((x+y+z)^2\le3(x^2+y^2+z^2)\), summing (9), and observing that the three collections on its right are disjoint subsets of \(\Omega_n\), we obtain
\[
\lVert A_{I Q}\rVert_F^2
\le3m^2\lVert P_{\Omega_n}A\rVert_F^2.                \tag{10}
\]
The missing block has rank at most two, and Cauchy--Schwarz applied to its singular values gives
\[
\lVert P_{W_n}\Delta\rVert_*
\le\sqrt2\lVert P_{W_n}\Delta\rVert_F
\le\sqrt6m\lVert P_{\Omega_n}\Delta\rVert_F.         \tag{11}
\]
Thus every wedge modulus, and hence its supremum over \(\Theta_n\), is a finite real number.

We now construct a uniformly valid relative-sign pilot.  Fix the same \(i_0\in C\) and \(t_0\in P\), set \(\operatorname{sign}(0)=1\), and define from observed coordinates
\[
\widehat p_i=\operatorname{sign}\!\left(\sum_{t\in P}X_{n,it}X_{n,i_0t}\right),
\qquad
\widehat q_t=\operatorname{sign}\!\left(\sum_{i\in C}X_{n,it}X_{n,it_0}\right).    \tag{12}
\]
All cells in (12) belong to \(\Omega_n\) by (1).  The observed-experiment definition, which follows from treatment support and the Gaussian-cell assumption, gives \(X_{n,it}=L_{n,it}+E_{n,it}\) on these coordinates.  For \(i\ne i_0\), multiplying the row score in (12) by \(p_ip_{i_0}\), and absorbing deterministic signs into the centered Gaussian errors, gives
\[
S_i=\sum_{t=1}^{m}(\ell_n+\varepsilon_t)(\ell_n+\delta_t)
=m\ell_n^2+G_i+H_i,                                    \tag{13}
\]
where the pairs \((\varepsilon_t,\delta_t)\) are independent over \(t\), the two coordinates in each pair are independent \(\mathcal N(0,\sigma^2)\), and
\[
G_i=\ell_n\sum_{t=1}^{m}(\varepsilon_t+\delta_t)
\sim\mathcal N(0,2m\ell_n^2\sigma^2),
\qquad H_i=\sum_{t=1}^{m}\varepsilon_t\delta_t.         \tag{14}
\]
For \(G\sim\mathcal N(0,v)\), exponential Markov inequality and minimization over \(a>0\) give
\[
\Pr(G\le-x)\le\inf_{a>0}\exp\{-ax+a^2v/2\}
=\exp\{-x^2/(2v)\}.                                   \tag{15}
\]
Applying (15) to (14) yields
\[
\Pr\!\left(G_i\le-\frac{m\ell_n^2}{2}\right)
\le\exp\!\left\{-\frac{m\ell_n^2}{16\sigma^2}\right\}.       \tag{16}
\]
For \(0<a<\sigma^{-2}\), conditioning on one Gaussian coordinate and integrating the other proves directly that
\[
\mathbb E\exp(-a\varepsilon_t\delta_t)
=(1-\sigma^4a^2)^{-1/2}.                                \tag{17}
\]
If \(0<a\le(2\sigma^2)^{-1}\), then \(-\tfrac12\log(1-\sigma^4a^2)\le\sigma^4a^2\).  Independence, (17), and exponential Markov inequality imply
\[
\Pr\!\left(H_i\le-\frac{m\ell_n^2}{2}\right)
\le\exp\!\left\{-\frac{ma\ell_n^2}{2}+m\sigma^4a^2\right\}.  \tag{18}
\]
Taking \(a=\min\{\ell_n^2/(4\sigma^4),(2\sigma^2)^{-1}\}\), and using \(\ell_n\ge s_->0\), gives
\[
\Pr\!\left(H_i\le-\frac{m\ell_n^2}{2}\right)\le e^{-mc_2},
\qquad
c_2=\min\!\left\{\frac{s_-^4}{16\sigma^4},\frac14\right\}>0. \tag{19}
\]
If \(S_i\le0\), at least one of the events in (16) and (19) occurs.  Thus, with
\[
c_1=\frac{s_-^2}{16\sigma^2},\qquad c_\star=\min\{c_1,c_2\}>0,
\]
we have
\[
\mathbb Q_{\vartheta,n}\{\widehat p_i\ne p_ip_{i_0}\}
\le e^{-mc_1}+e^{-mc_2}\qquad(i\ne i_0).               \tag{20}
\]
For \(i=i_0\), the row score is a sum of squares and is positive almost surely.  Interchanging rows and columns proves the same bound for \(\widehat q_t\ne q_tq_{t_0}\), while the reference-column score is positive almost surely.  Consequently, for
\[
\mathcal S_n=
\{\widehat p_i=p_ip_{i_0}\ \text{for every }i,
\ \widehat q_t=q_tq_{t_0}\ \text{for every }t\},
\]
the union bound, which requires no independence among the scores, yields
\[
\sup_{\vartheta_n\in\Theta_n}
\mathbb Q_{\vartheta,n}(\mathcal S_n^c)
\le2(n-1)(e^{-mc_1}+e^{-mc_2})
\le4ne^{-mc_\star}.                                    \tag{21}
\]

Suppose first that \(s_-<s_+\).  Let \(\Phi\) be the standard-normal distribution function and put \(z_\alpha=\Phi^{-1}(1-\alpha/4)>0\), using the coverage-level assumption.  Define
\[
\widehat\ell_n=\frac1{3m^2}\sum_{(i,t)\in\Omega_n}
X_{n,it}\widehat p_i\widehat q_t,
\qquad
\widehat L_{n,it}=\widehat\ell_n\widehat p_i\widehat q_t,
\qquad
R_n=\frac{z_\alpha}{\sqrt3}\sigma.                    \tag{22}
\]
On \(\mathcal S_n\), set \(c=p_{i_0}q_{t_0}\).  Then \(\widehat p_i\widehat q_t=cp_iq_t\), and (4), (22), and the independent Gaussian-cell assumption imply
\[
\widehat\ell_n=c(\ell_n+Z_n),
\qquad
Z_n=\frac1{3m^2}\sum_{(i,t)\in\Omega_n}p_iq_tE_{n,it}
\sim\mathcal N\!\left(0,\frac{\sigma^2}{3m^2}\right).           \tag{23}
\]
It follows on \(\mathcal S_n\) that
\[
P_{W_n}(\widehat L_n-L_n)=Z_np_Iq_Q^{\top},
\qquad
\lVert P_{W_n}(\widehat L_n-L_n)\rVert_*=m|Z_n|.       \tag{24}
\]
By the choice of \(z_\alpha\),
\[
\mathbb Q_{\vartheta,n}\{m|Z_n|>R_n\}=\alpha/2.       \tag{25}
\]
For all sufficiently large even \(n\), the right side of (21) is at most \(\alpha/2\).  The union bound applied to (21), (24), and (25), without any independence assertion between \(\mathcal S_n\) and \(Z_n\), proves
\[
\inf_{\vartheta_n\in\Theta_n}
\Pr_{\mathbb Q_{\vartheta,n},\Xi_n,\Zeta_n}
\!\left\{\lVert P_{W_n}(\widehat L_n-L_n)\rVert_*\le R_n\right\}
\ge1-\alpha.                                            \tag{26}
\]
Thus (22) belongs to \(\mathcal H_{\alpha,n}\) by the honest-procedures definition.  It is deterministic conditional on \(X_n\), and computing its row and column correlations, signed average, and outer-product completion takes \(O(n^2)\) arithmetic operations.  Since its radius is deterministic, the minimax-radius definition gives
\[
\rho_{\alpha,n}^{\star}\le C_\alpha\sigma,
\qquad C_\alpha:=\frac{z_\alpha}{\sqrt3}.             \tag{27}
\]

For the all-procedure lower bound, retain the fixed signs \(p_i=q_t=1\), and define
\[
a_\alpha=\frac{1-2\alpha}{2\sqrt3}>0,
\qquad
\delta_n=\frac{a_\alpha\sigma}{m},
\qquad
\ell_0=s_-,\qquad \ell_1=s_-+\delta_n.                \tag{28}
\]
Because \(s_-<s_+\), the condition \(\ell_1\le s_+\) holds for all sufficiently large \(n\).  Let
\[
L^{(j)}=\ell_j\mathbf1_n\mathbf1_n^{\top},qquad D^{(j)}=0,
\qquad j\in\{0,1\}.
\]
Equations (2)--(5) and treatment support verify that both tuples belong to the same \(\Theta_n\).  Their observed Gaussian mean vectors differ by \(\delta_n\) in each of the \(3m^2\) coordinates of \(\Omega_n\).

Let \(f_0,f_1\) be their observed Gaussian densities.  Completing the square in each coordinate gives the Hellinger affinity
\[
A:=\int\sqrt{f_0f_1}
=\exp\!\left\{-\frac{\lVert
\operatorname{vec}_{\Omega_n}(L^{(1)}-L^{(0)})\rVert_2^2}
{8\sigma^2}\right\}
=\exp\!\left\{-\frac{3a_\alpha^2}{8}\right\}.       \tag{29}
\]
Writing total variation as one half of the \(L^1\) distance and applying Cauchy--Schwarz gives
\[
\begin{aligned}
d_{\mathrm{TV}}(\mathbb Q_0,\mathbb Q_1)
&=\frac12\int|\sqrt{f_0}-\sqrt{f_1}|(\sqrt{f_0}+\sqrt{f_1})\\
&\le\frac12\{(2-2A)(2+2A)\}^{1/2}
=\sqrt{1-A^2}\\
&\le\frac{\sqrt3a_\alpha}{2}
=\frac{1-2\alpha}{4}.                                  \tag{30}
\end{aligned}
\]
The penultimate inequality follows from \(1-e^{-x}\le x\) for \(x\ge0\).  Adjoining the common independent laws of \(\Xi_n\) and \(\Zeta_n\) leaves total variation unchanged, because integration over that common product factor leaves the \(L^1\) distance unchanged.

Take an arbitrary \((\widehat L_n,R_n)\in\mathcal H_{\alpha,n}\), and on the common sample space define
\[
\mathcal A_j=
\{\lVert P_{W_n}(\widehat L_n-L^{(j)})\rVert_*\le R_n\},
\qquad j\in\{0,1\}.
\]
Uniform honesty gives \(\Pr_j(\mathcal A_j)\ge1-\alpha\).  Total variation gives \(\Pr_0(\mathcal A_1)\ge\Pr_1(\mathcal A_1)-d_{\mathrm{TV}}(\mathbb Q_0,\mathbb Q_1)\), so (30) implies
\[
\Pr_0(\mathcal A_0\cap\mathcal A_1)
\ge1-2\alpha-d_{\mathrm{TV}}(\mathbb Q_0,\mathbb Q_1)
\ge\frac34(1-2\alpha).                                 \tag{31}
\]
On this intersection, the triangle inequality for the nuclear norm gives
\[
2R_n\ge\lVert P_{W_n}(L^{(1)}-L^{(0)})\rVert_*
=\delta_n\lVert\mathbf1_I\mathbf1_Q^{\top}\rVert_*
=m\delta_n=a_\alpha\sigma.                             \tag{32}
\]
Therefore
\[
\sup_{\vartheta_n\in\Theta_n}
\mathbb E_{\mathbb Q_{\vartheta,n},\Xi_n,\Zeta_n}[R_n]
\ge\mathbb E_0[R_n]
\ge\frac{3a_\alpha(1-2\alpha)}8\sigma
=c_\alpha\sigma,
\qquad
c_\alpha:=\frac{3(1-2\alpha)^2}{16\sqrt3}>0.          \tag{33}
\]
Taking the infimum over all honest procedures in (33), and combining with (27), proves
\[
c_\alpha\sigma\le\rho_{\alpha,n}^{\star}\le C_\alpha\sigma,
\qquad \rho_{\alpha,n}^{\star}\asymp\sigma,          \tag{34}
\]
where both displayed constants depend only on \(\alpha\).  This lower bound varies only the positive factor level in a fixed-sign two-point subexperiment; the treatment-effect matrix is fixed at zero.  Thus it is an observed-law lower bound for the counterfactual radius, not an identification of a causal treatment effect by definition.

It remains to treat \(s_-=s_+\).  Write the known common value as \(s_0:=s_-=s_+>0\).  Use the relative signs in (12), and set
\[
\overline X_n^{\mathrm{sgn}}
=\frac1{3m^2}\sum_{(i,t)\in\Omega_n}X_{n,it}\widehat p_i\widehat q_t,
\qquad
\widehat c_n=\operatorname{sign}(\overline X_n^{\mathrm{sgn}}),
\]
\[
\widehat L_{n,it}=s_0\widehat c_n\widehat p_i\widehat q_t,
\qquad R_n=0.                                           \tag{35}
\]
On \(\mathcal S_n\), equation (23) becomes
\[
\overline X_n^{\mathrm{sgn}}=c(s_0+Z_n),
\qquad
Z_n\sim\mathcal N\!\left(0,\frac{\sigma^2}{3m^2}\right).       \tag{36}
\]
Hence \(\widehat c_n=c\) whenever \(\mathcal S_n\) occurs and \(Z_n>-s_0\), and on that event (35) equals \(L_n\) at every cell.  Applying (15) to (36) and using (21),
\[
\sup_{\vartheta_n\in\Theta_n}
\Pr_{\mathbb Q_{\vartheta,n},\Xi_n,\Zeta_n}
\{P_{W_n}\widehat L_n\ne P_{W_n}L_n\}
\le4ne^{-mc_\star}
+\exp\!\left\{-\frac{3m^2s_0^2}{2\sigma^2}\right\}.            \tag{37}
\]
For all sufficiently large even \(n\), the right side is at most \(\alpha\).  Therefore the deterministic \(O(n^2)\)-time pair in (35) belongs to \(\mathcal H_{\alpha,n}\).  Its radius is identically zero.  Since every radius in the honest-procedures definition is nonnegative, the minimax-radius definition yields
\[
0\le\rho_{\alpha,n}^{\star}\le0,
\qquad\text{so}\qquad \rho_{\alpha,n}^{\star}=0.      \tag{38}
\]

Finally, (8), (27), and finiteness from (11) imply in the nondegenerate case that
\[
0\le\frac{\rho_{\alpha,n}^{\star}}{\mathfrak R_n}
\le\frac{z_\alpha}{\sqrt{6n}}\longrightarrow0,         \tag{39}
\]
whereas the ratio is identically zero in the known-level case by (8) and (38).  Thus no fixed positive constant can make \(\rho_{\alpha,n}^{\star}\ge c\mathfrak R_n\) hold uniformly over all model settings and wedges admitted by the frozen assumptions.  Equations (4), (6), and (9) identify the obstruction: at \(\mu=1\), admissible singular-vector magnitudes are fixed and only coordinatewise signs remain, while the rank-manifold tangent space in the wedge-modulus definition still contains unrestricted continuous directions.
